\bookchapter{第五章：让手机更懂你——基础设置}

\chapterintro{花几分钟设置，体验提升巨大。这些设置就像给手机“调音”，让它更符合你的使用习惯。}

\booksection{一、声音与震动}

设个喜欢的铃声和通知声。灵活使用静音模式（开会）、勿扰模式（睡觉），世界瞬间清净。

\begin{itemize}
  \item \textbf{铃声}：选一个你能听到的、不尴尬的。
  \item \textbf{震动模式}：静音时靠震动提醒。
  \item \textbf{勿扰模式}：可以设置例外，比如允许通讯录来电。
  \item \textbf{音量分离}：媒体音量、铃声音量、闹钟音量可以分开调。
\end{itemize}

\booksection{二、显示与亮度}

\begin{itemize}
  \item \textbf{自动亮度}：让手机根据环境光自动调节，省心。
  \item \textbf{护眼/夜览模式}：过滤蓝光，晚上看手机舒服点。
  \item \textbf{深色模式}：省电（对OLED屏特别有效），且夜间不刺眼。很多App也支持深色模式。
  \item \textbf{屏幕刷新率}：部分安卓机可调，越高（如120Hz）滑动越丝滑。
  \item \textbf{自动锁屏}：设置15秒或30秒，安全又省电。
  \item \textbf{字体大小}：觉得字小？调大一点。
\end{itemize}

\booksection{三、壁纸、主题与桌面}

换张好看的壁纸（静态或动态），心情都变了。安卓品牌手机常有主题商店，能整套换图标和风格。

\textbf{桌面小组件}：天气、日程、音乐控制放桌面，一步直达。iOS和安卓都支持，长按桌面空白处添加。

\booksection{四、效率设置：让手机成为你的“外挂大脑”}

这部分是“一旦设置就回不去”的功能，每个都能让你的日常效率显著提升。

\booksubsection{专注 / 免打扰模式：夺回被通知夺走的注意力}

\begin{itemize}
  \item \textbf{核心思想}：不是"少看手机"，而是"在该专注时专注，该联系时再联系"。
  \item \textbf{使用建议}：\newline · 工作模式：仅允许"星标"联系人和必要 App 通过；\newline · 睡眠模式：自动开启，仅允许闹钟；\newline · 健身 / 阅读 / 会议：根据当前活动定制。
  \item \textbf{进阶玩法}：用自动化 + 位置 / 时间触发，让模式自动切换。例如：iOS 快捷指令可设置"到达公司即自动进入工作模式"。
\end{itemize}

\booksubsection{通知管理：让红点不再“夺命连环 call”}

\begin{itemize}
  \item \textbf{关闭营销通知}：90\% 的推送其实你并不需要。
  \item \textbf{按重要性排序}：把人际通讯类（电话、短信、即时通讯）置顶；新闻、购物、社交平台类可设为"静默推送"。
  \item \textbf{分组通知}：把同类通知折叠在一起，避免一次解锁被几十条消息淹没。
\end{itemize}

\booksubsection{快捷手势：让操作快人一步}

\begin{itemize}
  \item \textbf{截屏}：电源 + 音量下（iOS 全面屏）/电源 + 音量下（安卓通用）。
  \item \textbf{截长图 / 滚动截屏}：普通截图后点击"长截图"选项。
  \item \textbf{一键扫码}：控制中心 / 快捷设置里添加"扫一扫"图标。
  \item \textbf{轻点背面（iPhone）}：在 [设置] > [辅助功能] > [触控] > [轻点背面] 中绑定"截屏"或"锁屏"等动作。
  \item \textbf{语音唤醒}：Siri / 小爱同学 / 小艺 / 小布，喊一声就能查天气、设闹钟。
\end{itemize}

\booksubsection{输入法：最被低估的效率神器}

\begin{itemize}
  \item \textbf{选对输入法}：\newline · 搜狗 / 百度：词库大，流行语更新快；\newline · 讯飞：语音输入识别率最高，开会记笔记神器；\newline · Gboard（谷歌）：简洁、跨平台同步强、支持滑行输入；\newline · iOS 自带：隐私最好，与系统深度集成。
  \item \textbf{必备开关}：智能纠错、联想词、云词库同步、上滑输入数字符号。
  \item \textbf{隐私提示}：慎用"整句联想"或"个性化推荐"功能，可能上传输入内容。
\end{itemize}

\booksubsection{桌面与多任务：让工作流顺畅起来}

\begin{itemize}
  \item \textbf{App 分屏 / 画中画}：看视频时回消息，上网时记笔记。
  \item \textbf{小窗 / 悬浮窗}：微信 / 抖音支持"小窗模式"，回复完即收起。
  \item \textbf{常用 App 置顶}：Dock 栏（屏幕底部）只放最高频的 4 个，减少查找时间。
\end{itemize}

\begin{tcbinfo}{"上工第一天"检查清单}
拿到新手机或大版本更新后，建议花 10 分钟过一遍这些设置：\newline (1) 锁屏密码 + 生物识别 (2) 通知分组 + 关闭营销推送 (3) 字体大小 / 刷新率 / 深色模式\newline (4) 默认输入法 / 浏览器 / 邮件 (5) 紧急联系人 / 医疗卡 (6) 备份设置

\end{tcbinfo}

\booksection{五、辅助功能与无障碍：科技向善的力量}

别被“辅助功能”这个名字吓到，这些功能人人能用。它们最初是为残障人士设计的，但对普通人也很有用——临时手受伤抱娃、地铁上耳机坏了听不清、戴老花镜看小字……都能救场。说到底，无障碍设计体现了\textbf{科技向善}的理念：让每个人都能平等地享受数字生活。

\booksubsection{视障用户的“眼睛”——屏幕阅读器全攻略}

屏幕阅读器是视障用户使用手机的核心工具。它会把屏幕上\textbf{所有可视内容朗读出来}，用户通过手势和语音操作手机。下面分别介绍 iOS “旁白”和安卓 TalkBack 的使用方法。

\begin{tcbinfo}{iOS 旁白（VoiceOver）}
\textbf{开启路径}：[设置] > [辅助功能] > [旁白]，打开开关。也可三击侧边键（或主屏幕键）快速开关。

\textbf{基本手势}：

\begin{enumerate}
  \item 单指滑动：逐元素浏览屏幕内容；
  \item 双指轻扫：上下/左右整页滚动；
  \item 单指双击：激活当前焦点元素（相当于"点击"）；
  \item 三指轻扫：切换页面；
  \item 单指三击：放大或选择（与普通三击不同）。
\end{enumerate}

\textbf{个性化设置}：[旁白] > [语音]，可调节\textbf{朗读语速}（默认约 50\%，可调到 70-80\% 听感更流畅）、\textbf{音调}、切换不同语音包。iOS 自带“增强语音”和“语音控制器”，前者能减少呼吸声和停顿，听感更自然。

\end{tcbinfo}

\begin{tcbinfo}{安卓 TalkBack}
\textbf{开启路径}：[设置] > [辅助功能] > [TalkBack]，打开开关。或长按音量上下键 3 秒快速启用（部分品牌支持）。

\textbf{基本手势}：

\begin{enumerate}
  \item 单指左右滑动：切换焦点元素；
  \item 单指双击：激活元素；
  \item 双指上下/左右滑：滚动；
  \item 三指单击：进入阅读模式（连续朗读）；
  \item 四指单击：暂停/继续朗读。
\end{enumerate}

\textbf{个性化设置}：[TalkBack] > [文字转语音设置]，可调整\textbf{语速}、\textbf{音调}、切换\textbf{语音引擎}（建议安装 Google 文字转语音引擎，朗读更自然）。

\end{tcbinfo}

\booksubsection{让小字变大、让模糊变清晰——视觉调整}

\booksubsection{放大手势}

\begin{itemize}
  \item \textbf{iOS}：[设置] > [辅助功能] > [缩放]，开启后\textbf{三指轻点屏幕}即可放大。默认全屏缩放，可选窗口缩放（局部放大窗）。\textbf{双指开合可调节倍数（1.5x 至 15x）}，三指双击锁定当前缩放比例。
  \item \textbf{安卓}：[设置] > [辅助功能] > [放大]，三指点击放大，再三指点击退出；或直接\textbf{三指点击屏幕三下}启用临时放大（最多 8 倍）。
\end{itemize}

\booksubsection{高对比度文字 / 粗体文本}

\begin{itemize}
  \item \textbf{iOS}：[设置] > [辅助功能] > [显示与文字大小]，开启"粗体文本"、"增强对比度"、"降低透明度"等开关。"智能反转"可在浅色背景下让文字保持黑色，避免白底白字。
  \item \textbf{安卓}：[设置] > [显示] > [字体大小]，可调节"小、普通、大、超大"四档。华为/小米额外提供"粗体"和"高对比度文本"开关。
\end{itemize}

\booksubsection{色彩滤镜（色盲友好）}

\begin{itemize}
  \item \textbf{iOS}：[设置] > [辅助功能] > [显示与文字大小] > [色彩滤镜]，提供\textbf{红/绿/蓝/黄四种色盲模式}，以及"灰度"和"反转颜色"选项。建议先做一次色觉测试（系统自带）再选对应滤镜。
  \item \textbf{安卓}：[设置] > [辅助功能] > [颜色校正]，含\textbf{三种色盲模式}+ "反转颜色"开关。
\end{itemize}

\booksubsection{AI 图像描述：让 AI 当你的“眼睛”}

当屏幕阅读器遇到图片时，传统做法是跳过或只读出“图片”两个字。但现在 AI 可以\textbf{用一句话描述图片内容}，对识别验证码、读图、识别菜单非常有帮助。

\begin{tcbinfo}{各系统 AI 图像描述支持情况}
\begin{itemize}
  \item \textbf{iOS 15+}：在 [设置] > [辅助功能] > [旁白] > [语音描述] 中开启"图像描述"。系统会调用\textbf{本地 AI 模型}自动生成图片文字描述（需 iPhone 12 及以上）。
  \item \textbf{新版 TalkBack 应用}：Google TalkBack 在较新版本中集成\textbf{图片问答}功能，可以问"这张图里有什么"，AI 实时回答（不严格依赖安卓系统版本，更新 TalkBack 应用即可体验）。
  \item \textbf{第三方 App}："Seeing AI"（iOS）、"Lookout"（安卓）、"小艺帮帮"（华为）都是优秀的 AI 视觉助手，可识别物品、读文件、识别人脸表情。
\end{itemize}

\end{tcbinfo}

\begin{tcbpractice}{使用场景}
试着对着一张街拍照片问“这张图里有几个人？他们在做什么？”，或者打开外卖 App，让 AI 读出菜单图里的菜品名——\textbf{这是真正的“科技让生活更平等”}。

\end{tcbpractice}

\booksubsection{听障用户的“耳朵”——实时字幕全攻略}

实时字幕（Live Captions）是\textbf{系统级的语音转文字}功能，听障用户在嘈杂环境、图书馆、外语场景下都特别有用。

\begin{tcbinfo}{实时字幕开启与优化}
\textbf{iOS}：[设置] > [辅助功能] > [实时字幕]，打开开关。系统会检测设备所有音频（通话、视频、播客、FaceTime、语音备忘录），自动显示字幕在屏幕上方。

\textbf{安卓 12+}：按\textbf{音量键}点击弹出框中的“实时字幕”图标；或在 [设置] > [辅助功能] > [实时字幕] 中开启。

\textbf{准确率优化技巧}：

\begin{enumerate}
  \item 尽量在\textbf{安静环境}下使用，远离空调、机器轰鸣；
  \item 说话者\textbf{咬字清晰、语速适中}时识别率最高；
  \item 多人对话会降低准确率，可使用"只显示一人"模式；
  \item 支持中文普通话、粤语、英语等主流语种，\textbf{小语种识别率较低}；
  \item 说话中途有背景音乐或人声干扰，识别会断断续续。
\end{enumerate}

\end{tcbinfo}

\booksubsection{通话场景双通道：语音转文字 + 文字转语音}

当用户\textbf{既听不到也说不了话}时（如突发失声、喉部手术恢复期、严重口吃），手机还能当“传话员”。

\begin{tcbinfo}{通话实时转译（TTY / RTT）}
\textbf{iOS}：[设置] > [辅助功能] > [RTT/TTY]，开启后通话时屏幕上会显示键盘，对方的语音会被实时打成文字，\textbf{你打的字也会被念出来}给对方听。

\textbf{安卓}：[设置] > [辅助功能] > [实时转录] + [声音通知]。“实时转录”将对方说话实时显示为文字；“声音通知”在不便语音时将你的输入转为语音播放。

\textbf{兼容性}：

\begin{itemize}
  \item \textbf{RTT（Real-Time Text）}：双方都需支持，文字边打边显示，无延迟，主流运营商已支持；
  \item \textbf{TTY（电传打字）}：需专用 TTY 设备，已逐步被 RTT 取代；
  \item 微信/QQ 的\textbf{语音转文字功能}也是替代方案：长按语音消息 "转文字"。
\end{itemize}

\end{tcbinfo}

\booksubsection{面向所有人的“小但实用”功能}

\begin{itemize}
  \item \textbf{轻点背面}（iOS 14+）：[设置] > [辅助功能] > [触控] > [轻点背面]，可绑定\textbf{双击/三击}触发截图、锁屏、控制中心等动作。双手拿快递时特别好用。
  \item \textbf{辅助触控 / 悬浮球}：iOS [辅助功能] > [触控] > [辅助触控]，安卓 [辅助功能] > [悬浮导航]。虚拟按钮替代手势或物理按键，按键不灵或单手操作时救场。
  \item \textbf{切换控制}：通过外部开关、单次点击、头部动作控制手机，\textbf{为肢体障碍用户设计}。iOS [辅助功能] > [切换控制]，安卓 [辅助功能] > [开关控制]。
  \item \textbf{声音识别}（iOS）：[辅助功能] > [声音识别]，开启后手机会"听"门铃、婴儿哭声、火警警报、警笛等声音，\textbf{识别到就震动+屏幕通知}，对听障用户和睡眠期间非常有价值。
  \item \textbf{LED 闪烁提醒}：iOS [辅助功能] > [音频/视觉] > [LED 闪烁以示提醒]，安卓 [设置] > [通知] > [闪光灯提醒]。来电和通知时闪光灯闪烁，\textbf{嘈杂环境无声也能感知}。
  \item \textbf{单声道音频}：[设置] > [辅助功能] > [音频] > [单声道音频]，将左右声道合并。单耳听力受损的用户不会丢失另一半声音信息。
\end{itemize}

\begin{tcbinfo}{苹果的无障碍设计哲学}
苹果在无障碍领域投入巨大，每年都会更新和新增辅助功能。苹果官网有专门的无障碍页面，详细介绍各项功能的使用方法。iOS 的无障碍功能通常比安卓更完善、体验更流畅，这也是很多残障用户选择苹果的重要原因。

\end{tcbinfo}

\begin{tcbpractice}{花30秒体验一下}
找一段空闲时间，做以下体验：

\begin{enumerate}
  \item 开启\textbf{旁白 / TalkBack}，试着不靠眼睛用三指滑动浏览微信；
  \item 打开\textbf{实时字幕}，对着手机说话看识别准确率；
  \item 用\textbf{三指点击}放大一段文字到 200\%，体验视障用户的日常；
  \item 关掉之前所有功能，思考"如果我永远需要这些功能，我会希望它怎么设计"。
\end{enumerate}

\end{tcbpractice}
